% Options for packages loaded elsewhere
\PassOptionsToPackage{unicode}{hyperref}
\PassOptionsToPackage{hyphens}{url}
\documentclass[
  11pt,
]{article}
\usepackage{xcolor}
\usepackage[margin=1in]{geometry}
\usepackage{amsmath,amssymb}
\setcounter{secnumdepth}{-\maxdimen} % remove section numbering
\usepackage{iftex}
\ifPDFTeX
  \usepackage[T1]{fontenc}
  \usepackage[utf8]{inputenc}
  \usepackage{textcomp} % provide euro and other symbols
\else % if luatex or xetex
  \usepackage{unicode-math} % this also loads fontspec
  \defaultfontfeatures{Scale=MatchLowercase}
  \defaultfontfeatures[\rmfamily]{Ligatures=TeX,Scale=1}
\fi
\usepackage{lmodern}
\ifPDFTeX\else
  % xetex/luatex font selection
\fi
% Use upquote if available, for straight quotes in verbatim environments
\IfFileExists{upquote.sty}{\usepackage{upquote}}{}
\IfFileExists{microtype.sty}{% use microtype if available
  \usepackage[]{microtype}
  \UseMicrotypeSet[protrusion]{basicmath} % disable protrusion for tt fonts
}{}
\makeatletter
\@ifundefined{KOMAClassName}{% if non-KOMA class
  \IfFileExists{parskip.sty}{%
    \usepackage{parskip}
  }{% else
    \setlength{\parindent}{0pt}
    \setlength{\parskip}{6pt plus 2pt minus 1pt}}
}{% if KOMA class
  \KOMAoptions{parskip=half}}
\makeatother
\usepackage{graphicx}
\makeatletter
\newsavebox\pandoc@box
\newcommand*\pandocbounded[1]{% scales image to fit in text height/width
  \sbox\pandoc@box{#1}%
  \Gscale@div\@tempa{\textheight}{\dimexpr\ht\pandoc@box+\dp\pandoc@box\relax}%
  \Gscale@div\@tempb{\linewidth}{\wd\pandoc@box}%
  \ifdim\@tempb\p@<\@tempa\p@\let\@tempa\@tempb\fi% select the smaller of both
  \ifdim\@tempa\p@<\p@\scalebox{\@tempa}{\usebox\pandoc@box}%
  \else\usebox{\pandoc@box}%
  \fi%
}
% Set default figure placement to htbp
\def\fps@figure{htbp}
\makeatother
\setlength{\emergencystretch}{3em} % prevent overfull lines
\providecommand{\tightlist}{%
  \setlength{\itemsep}{0pt}\setlength{\parskip}{0pt}}
\usepackage{bookmark}
\IfFileExists{xurl.sty}{\usepackage{xurl}}{} % add URL line breaks if available
\urlstyle{same}
\hypersetup{
  pdftitle={Introduction à l'économétrie},
  hidelinks,
  pdfcreator={LaTeX via pandoc}}

\title{Introduction à l'économétrie}
\author{}
\date{\vspace{-2.5em}Printemps 2026}

\begin{document}
\maketitle

\section{Organisation du cours}\label{organisation-du-cours}

\textbf{Enseignant(s) :} \emph{À compléter}\\
\textbf{Email :} \emph{À compléter}\\
\textbf{Organisation :} 12 séances.

\section{Description du cours}\label{description-du-cours}

Ce cours propose une introduction aux méthodes fondamentales de
l'économétrie et à leur mise en œuvre avec \textbf{R}. Il vise à donner
aux étudiants les outils nécessaires pour analyser des données
économiques, estimer et interpréter des modèles de régression, réaliser
des tests statistiques et comprendre les principales méthodes
d'identification causale utilisées en économie empirique. Une attention
particulière est portée à la reproductibilité des analyses et à
l'interprétation économique des résultats.

Chaque séance combine des éléments théoriques, des illustrations
empiriques et des applications sous \textbf{R}.

\section{Logiciels}\label{logiciels}

\begin{itemize}
\tightlist
\item
  \textbf{R} et \textbf{RStudio}. Les étudiants devront installer les
  dernières versions de \textbf{R} et \textbf{RStudio} avant le premier
  cours.
\item
  Une connaissance élémentaire de \textbf{R Markdown} est recommandée.
\end{itemize}

\section{Prérequis}\label{pruxe9requis}

Connaissances de base en statistiques descriptives et en économie.

\section{Supports pédagogiques}\label{supports-puxe9dagogiques}

Les diapositives de cours, les jeux de données et les scripts \textbf{R}
seront mis à disposition au fur et à mesure du semestre.

Les principaux ouvrages de référence sont :

\begin{itemize}
\tightlist
\item
  \emph{Introduction to Econometrics with R}
\item
  \emph{Mastering Metrics} (Angrist \& Pischke)
\item
  \emph{ModernDive}
\item
  \emph{R for Data Science}
\end{itemize}

Les lectures indiquées en \textbf{gras} sont \textbf{fortement
recommandées}.

\section{Plan du cours}\label{plan-du-cours}

\subsection{Séance 1 : Introduction}\label{suxe9ance-1-introduction}

\textbf{Contenu}

\begin{itemize}
\tightlist
\item
  Qu'est-ce que l'économétrie ?
\item
  Corrélation et causalité
\item
  Introduction à R
\item
  La « Credibility Revolution » en économie empirique
\end{itemize}

\textbf{Lectures recommandées}

\begin{itemize}
\tightlist
\item
  \textbf{Sections 1.1 et 1.2, \emph{Introduction to Econometrics with
  R}}
\item
  \textbf{Introduction, \emph{Mastering Metrics}}
\item
  \textbf{Angrist \& Pischke (2010), \emph{The Credibility Revolution in
  Empirical Economics}}
\end{itemize}

\begin{center}\rule{0.5\linewidth}{0.5pt}\end{center}

\subsection{Séances 2 et 3 : Manipulation, visualisation et description
des
données}\label{suxe9ances-2-et-3-manipulation-visualisation-et-description-des-donnuxe9es}

\textbf{Contenu}

\begin{itemize}
\tightlist
\item
  Manipulation des données avec le \emph{tidyverse}
\item
  Nettoyage des données
\item
  Statistiques descriptives
\item
  Visualisation graphique des données
\end{itemize}

\textbf{Lectures recommandées}

\begin{itemize}
\tightlist
\item
  \textbf{Chapitre « Data Transformation », \emph{R for Data Science}}
\item
  \textbf{Schwabish (2014), \emph{An Economist's Guide to Visualizing
  Data}}
\item
  Chapitres 2 et 3 de \emph{ModernDive}
\end{itemize}

\begin{center}\rule{0.5\linewidth}{0.5pt}\end{center}

\subsection{Séance 4 : Régression linéaire
simple}\label{suxe9ance-4-ruxe9gression-linuxe9aire-simple}

\textbf{Contenu}

\begin{itemize}
\tightlist
\item
  Le modèle de régression linéaire simple
\item
  Estimation par les moindres carrés ordinaires (MCO)
\item
  Interprétation des coefficients
\item
  Qualité de l'ajustement
\end{itemize}

\textbf{Lectures recommandées}

\begin{itemize}
\tightlist
\item
  \textbf{Chapitre 3, \emph{Introduction to Econometrics with R}}
\item
  Chapitre 5 de \emph{ModernDive}
\end{itemize}

\begin{center}\rule{0.5\linewidth}{0.5pt}\end{center}

\subsection{Séance 5 : Introduction à la
causalité}\label{suxe9ance-5-introduction-uxe0-la-causalituxe9}

\textbf{Contenu}

\begin{itemize}
\tightlist
\item
  Résultats potentiels
\item
  Raisonnement contrefactuel
\item
  Effet moyen d'un traitement
\item
  Identification des effets causaux
\end{itemize}

\textbf{Lectures recommandées}

\begin{itemize}
\tightlist
\item
  \textbf{Chapitre 7, \emph{Introduction to Econometrics with R}}
\item
  \textbf{Chapitre 1, \emph{Mastering Metrics}}
\item
  Cunningham, \emph{Causal Inference: The Mixtape}
\item
  Morgan \& Winship (2015)
\end{itemize}

\begin{center}\rule{0.5\linewidth}{0.5pt}\end{center}

\subsection{Séance 6 : Régression linéaire
multiple}\label{suxe9ance-6-ruxe9gression-linuxe9aire-multiple}

\textbf{Contenu}

\begin{itemize}
\tightlist
\item
  Régression multiple
\item
  Variables de contrôle
\item
  Biais de variable omise
\item
  Interprétation des coefficients
\end{itemize}

\textbf{Lectures recommandées}

\begin{itemize}
\tightlist
\item
  \textbf{Chapitre 4, \emph{Introduction to Econometrics with R}}
\item
  \textbf{Chapitre 2, \emph{Mastering Metrics}}
\item
  Chapitre 6 de \emph{ModernDive}
\end{itemize}

\begin{center}\rule{0.5\linewidth}{0.5pt}\end{center}

\subsection{Séance 7 :
Échantillonnage}\label{suxe9ance-7-uxe9chantillonnage}

\textbf{Contenu}

\begin{itemize}
\tightlist
\item
  Échantillonnage aléatoire
\item
  Distributions d'échantillonnage
\item
  Introduction au bootstrap
\end{itemize}

\textbf{Lectures recommandées}

\begin{itemize}
\tightlist
\item
  \textbf{Chapitre 7 de \emph{ModernDive}}
\end{itemize}

\begin{center}\rule{0.5\linewidth}{0.5pt}\end{center}

\subsection{Séance 8 : Intervalles de confiance et tests
d'hypothèses}\label{suxe9ance-8-intervalles-de-confiance-et-tests-dhypothuxe8ses}

\textbf{Contenu}

\begin{itemize}
\tightlist
\item
  Inférence statistique
\item
  Intervalles de confiance
\item
  Tests d'hypothèses
\item
  Valeurs \emph{p}
\end{itemize}

\textbf{Lectures recommandées}

\begin{itemize}
\tightlist
\item
  \textbf{Chapitres 8 et 9 de \emph{ModernDive}}
\end{itemize}

\begin{center}\rule{0.5\linewidth}{0.5pt}\end{center}

\subsection{Séance 9 : Inférence en
régression}\label{suxe9ance-9-infuxe9rence-en-ruxe9gression}

\textbf{Contenu}

\begin{itemize}
\tightlist
\item
  Erreurs standards
\item
  Significativité statistique
\item
  Intervalles de confiance
\item
  Tests sur les coefficients de régression
\end{itemize}

\textbf{Lectures recommandées}

\begin{itemize}
\tightlist
\item
  \textbf{Chapitre 6, \emph{Introduction to Econometrics with R}}
\item
  \textbf{Chapitre 10 de \emph{ModernDive}}
\end{itemize}

\begin{center}\rule{0.5\linewidth}{0.5pt}\end{center}

\subsection{\texorpdfstring{Séance 10 : La méthode des doubles
différences
(\emph{Differences-in-Differences})}{Séance 10 : La méthode des doubles différences (Differences-in-Differences)}}\label{suxe9ance-10-la-muxe9thode-des-doubles-diffuxe9rences-differences-in-differences}

\textbf{Contenu}

\begin{itemize}
\tightlist
\item
  Estimateur en doubles différences
\item
  Hypothèse de tendances parallèles
\item
  Effets fixes
\item
  Applications empiriques
\end{itemize}

\textbf{Lectures recommandées}

\begin{itemize}
\tightlist
\item
  \textbf{Chapitre 5, \emph{Mastering Metrics}}
\item
  \textbf{Card \& Krueger (1994), « Minimum Wages and Employment: A Case
  Study of the Fast-Food Industry in New Jersey and Pennsylvania »}
\end{itemize}

\begin{center}\rule{0.5\linewidth}{0.5pt}\end{center}

\subsection{Séance 11 : Régression sur
discontinuité}\label{suxe9ance-11-ruxe9gression-sur-discontinuituxe9}

\textbf{Contenu}

\begin{itemize}
\tightlist
\item
  Discontinuité franche et floue
\item
  Hypothèses d'identification
\item
  Estimation et interprétation
\item
  Applications empiriques
\end{itemize}

\textbf{Lectures recommandées}

\begin{itemize}
\tightlist
\item
  \textbf{Chapitre 4, \emph{Mastering Metrics}}
\item
  \textbf{Les diapositives et vidéos d'Andrew Heiss}
\item
  Chris Walters (2020), \emph{Regression Discontinuity Designs}
\item
  \textbf{Carpenter \& Dobkin (2009), « The Effect of Alcohol
  Consumption on Mortality »}
\end{itemize}

\textbf{Lectures complémentaires}

\begin{itemize}
\tightlist
\item
  Imbens \& Lemieux (2008)
\item
  Lee \& Lemieux (2010)
\end{itemize}

\begin{center}\rule{0.5\linewidth}{0.5pt}\end{center}

\subsection{Séance 12 : Séance de
révision}\label{suxe9ance-12-suxe9ance-de-ruxe9vision}

\textbf{Contenu}

\begin{itemize}
\tightlist
\item
  Synthèse des notions abordées
\item
  Exercices de révision
\item
  Questions et discussion
\end{itemize}

\section{Références}\label{ruxe9fuxe9rences}

\subsection{Ouvrages principaux}\label{ouvrages-principaux}

\begin{itemize}
\tightlist
\item
  \textbf{Angrist, J. D. \& Pischke, J.-S. (2015). \emph{Mastering
  Metrics}.}
\item
  \textbf{Introduction to Econometrics with R.}
\item
  \textbf{ModernDive.}
\item
  \textbf{Wickham, H. \& Grolemund, G. \emph{R for Data Science}.}
\end{itemize}

\subsection{Articles de référence}\label{articles-de-ruxe9fuxe9rence}

\begin{itemize}
\tightlist
\item
  Angrist, J. D. \& Pischke, J.-S. (2010). \emph{The Credibility
  Revolution in Empirical Economics.}
\item
  Schwabish, J. (2014). \emph{An Economist's Guide to Visualizing Data.}
\item
  Card, D. \& Krueger, A. (1994). \emph{Minimum Wages and Employment: A
  Case Study of the Fast-Food Industry in New Jersey and Pennsylvania.}
\item
  Carpenter, C. \& Dobkin, C. (2009). \emph{The Effect of Alcohol
  Consumption on Mortality: Regression Discontinuity Evidence from the
  Minimum Drinking Age.}
\end{itemize}

\subsection{Lectures complémentaires}\label{lectures-compluxe9mentaires}

\begin{itemize}
\tightlist
\item
  Cunningham, S. \emph{Causal Inference: The Mixtape.}
\item
  Morgan, S. L. \& Winship, C. \emph{Counterfactuals and Causal
  Inference.}
\item
  Imbens, G. \& Lemieux, T. (2008).
\item
  Lee, D. \& Lemieux, T. (2010).
\end{itemize}

\end{document}
